% Completed frozen F0 results and explicitly prospective experiment settings.
\subsection{Complete PAMF: CRIC with DURA}
\label{sec:cric-dura}
The user-designated complete PAMF system applies the frozen DURA v1.0 restorer to each received IQ window, followed by CRIC and its original jointly trained cosine classifier. This includes clean observations; the system does not use the true channel scene to decide whether restoration is applied. CRIC denotes Cross-Receiver Identity Contrastive Learning, the original E50 model with 33 labeled-only epochs and 17 joint labeled/unlabeled epochs. It is distinct from the later SARD/T1 schedule with 13 labeled-only and 37 joint epochs.

CRIC retains three SIG\_TIME paths, time/frequency/PA residual fusion, reference-response features, centered and unit-aware Sinc optimization, R-Drop, label smoothing 0.01, and its original EMA-based unlabeled selection. Cross-RX contrast uses actual labeled features from shallow and intermediate layers, with same-TX/different-physical-RX positives and different-TX negatives. Same-TX/same-RX pairs are excluded. Temperature is 0.2; weight is zero for five epochs and rises linearly over ten epochs to 0.1. The two existing R-Drop forwards supply both levels; contrast adds no network forward. The original CRIC has no satellite training branch. It is preserved as a complete frozen recipe, rather than silently replacing its EMA/gates with ordinary PL or adding SARD augmentation.

DURA uses two learned-prior stages and two conjugate-gradient updates per stage. Its 266951 unique parameters predict a coarse restoration and a current-window real-linear action, with a main and conjugate-image response. The fixed runtime uses residual reuse and CUDA Graphs in FP32 with TF32 disabled. Its 27 FFT/IFFT calls describe G only. Restorer training uses source IQ and simulated pairs without TX labels, identity-model weights/outputs, or target observations. All training ancestors are subject to the same source contract.

% HISTORICAL old no-As 33L17U F, NOT final T1-E50.
\begin{table*}[t]
\centering\scriptsize\setlength{\tabcolsep}{3pt}
\caption{Historical F0 accuracy (\%): old no-As 33L/17U identifier, not final T1; frozen complete recipes, mean $\pm$ sample SD over three F seeds. All use the same seven target RX and 168000 packets per view. CRIC+DURA conditions on one frozen G. S/U denote suburban/urban; S6 is the per-seed mean of six satellite views.}
\label{tab:cric-full-f0}
\begin{tabular}{llrrrrrrrr}
\toprule
Method & Training & Clean & S-high & S-mid & S-low & U-high & U-mid & U-low & S6 \\
\midrule
RIEI CE & Supervised & $66.50\!\pm\!1.37$ & $61.66\!\pm\!1.02$ & $59.09\!\pm\!1.03$ & $57.22\!\pm\!0.89$ & $58.47\!\pm\!1.00$ & $45.61\!\pm\!0.79$ & $36.68\!\pm\!0.63$ & $53.12\!\pm\!0.89$ \\
RIEI Full & Supervised & $54.53\!\pm\!8.31$ & $55.21\!\pm\!8.39$ & $53.38\!\pm\!7.62$ & $52.02\!\pm\!6.86$ & $52.03\!\pm\!7.82$ & $40.47\!\pm\!4.87$ & $32.97\!\pm\!2.93$ & $47.68\!\pm\!6.42$ \\
DRIFT CE & Supervised & $68.49\!\pm\!2.50$ & $37.07\!\pm\!5.55$ & $36.20\!\pm\!5.07$ & $35.34\!\pm\!4.64$ & $36.06\!\pm\!5.07$ & $31.87\!\pm\!3.42$ & $28.59\!\pm\!2.30$ & $34.19\!\pm\!4.33$ \\
DRIFT Full & Supervised & $67.33\!\pm\!5.91$ & $29.08\!\pm\!4.48$ & $27.96\!\pm\!3.97$ & $27.04\!\pm\!3.77$ & $28.42\!\pm\!4.14$ & $25.37\!\pm\!2.62$ & $23.10\!\pm\!2.01$ & $26.83\!\pm\!3.36$ \\
IB-RFF CE & Supervised & $63.30\!\pm\!2.14$ & $18.27\!\pm\!3.02$ & $18.09\!\pm\!2.80$ & $17.97\!\pm\!2.66$ & $18.24\!\pm\!2.67$ & $17.93\!\pm\!1.71$ & $17.74\!\pm\!1.10$ & $18.04\!\pm\!2.31$ \\
IB-RFF Full & Supervised & $63.77\!\pm\!2.18$ & $15.22\!\pm\!6.07$ & $15.41\!\pm\!5.61$ & $15.52\!\pm\!5.32$ & $15.43\!\pm\!5.59$ & $16.35\!\pm\!3.88$ & $16.74\!\pm\!2.83$ & $15.78\!\pm\!4.87$ \\
SDG CE & Supervised & $73.73\!\pm\!0.91$ & $36.77\!\pm\!6.30$ & $36.63\!\pm\!6.03$ & $36.19\!\pm\!5.90$ & $36.37\!\pm\!5.82$ & $33.97\!\pm\!4.07$ & $30.94\!\pm\!2.90$ & $35.15\!\pm\!5.17$ \\
SDG Full & Supervised & $71.39\!\pm\!1.84$ & $65.95\!\pm\!2.85$ & $63.97\!\pm\!2.81$ & $61.99\!\pm\!2.77$ & $63.41\!\pm\!2.62$ & $51.54\!\pm\!1.82$ & $42.00\!\pm\!1.10$ & $58.14\!\pm\!2.32$ \\
RIEI CE & PL & $66.45\!\pm\!1.37$ & $61.56\!\pm\!1.10$ & $58.98\!\pm\!1.08$ & $57.11\!\pm\!0.97$ & $58.37\!\pm\!1.11$ & $45.45\!\pm\!0.86$ & $36.52\!\pm\!0.70$ & $53.00\!\pm\!0.97$ \\
RIEI Full & PL & $53.53\!\pm\!6.48$ & $54.92\!\pm\!7.39$ & $53.18\!\pm\!6.73$ & $51.93\!\pm\!6.17$ & $51.73\!\pm\!6.80$ & $40.34\!\pm\!4.28$ & $32.91\!\pm\!2.60$ & $47.50\!\pm\!5.66$ \\
DRIFT CE & PL & $70.50\!\pm\!3.78$ & $32.80\!\pm\!8.75$ & $32.58\!\pm\!8.08$ & $32.24\!\pm\!7.59$ & $32.17\!\pm\!8.16$ & $29.50\!\pm\!5.52$ & $27.14\!\pm\!3.84$ & $31.07\!\pm\!6.99$ \\
DRIFT Full & PL & $60.97\!\pm\!11.26$ & $29.57\!\pm\!3.69$ & $28.50\!\pm\!3.61$ & $27.58\!\pm\!3.60$ & $28.89\!\pm\!3.30$ & $25.70\!\pm\!2.55$ & $23.44\!\pm\!1.90$ & $27.28\!\pm\!3.11$ \\
IB-RFF CE & PL & $64.78\!\pm\!3.20$ & $17.77\!\pm\!0.98$ & $17.92\!\pm\!0.87$ & $17.94\!\pm\!0.90$ & $17.74\!\pm\!0.92$ & $17.73\!\pm\!0.76$ & $17.65\!\pm\!0.52$ & $17.79\!\pm\!0.81$ \\
IB-RFF Full & PL & $66.08\!\pm\!0.99$ & $18.55\!\pm\!5.53$ & $18.48\!\pm\!5.11$ & $18.44\!\pm\!4.74$ & $18.52\!\pm\!5.03$ & $18.46\!\pm\!3.26$ & $18.16\!\pm\!2.14$ & $18.44\!\pm\!4.30$ \\
SDG CE & PL & $72.01\!\pm\!2.78$ & $46.01\!\pm\!6.46$ & $44.89\!\pm\!6.37$ & $43.74\!\pm\!6.37$ & $44.80\!\pm\!5.94$ & $38.88\!\pm\!4.16$ & $33.79\!\pm\!2.96$ & $42.02\!\pm\!5.37$ \\
SDG Full & PL & $72.10\!\pm\!0.41$ & $65.81\!\pm\!1.64$ & $63.86\!\pm\!1.60$ & $61.99\!\pm\!1.64$ & $63.23\!\pm\!1.56$ & $51.70\!\pm\!1.46$ & $42.32\!\pm\!1.33$ & $58.15\!\pm\!1.53$ \\
\midrule
CRIC raw & Native SSL & $85.11\!\pm\!0.38$ & $49.26\!\pm\!7.38$ & $47.74\!\pm\!6.77$ & $46.09\!\pm\!6.12$ & $47.49\!\pm\!6.75$ & $39.69\!\pm\!4.52$ & $33.65\!\pm\!2.93$ & $43.99\!\pm\!5.74$ \\
CRIC$+$DURA & Native SSL & $85.14\!\pm\!0.41$ & $82.80\!\pm\!0.35$ & $81.45\!\pm\!0.28$ & $79.85\!\pm\!0.25$ & $81.62\!\pm\!0.32$ & $74.90\!\pm\!0.53$ & $68.13\!\pm\!0.66$ & $78.12\!\pm\!0.34$ \\
\bottomrule
\end{tabular}
\end{table*}


\subsubsection{Completed F0 Comparison}
Table~\ref{tab:cric-full-f0} combines existing immutable scores; it entails no new training or target-truth access. Its 18 method/recipe readouts each have three F seeds and seven views, giving 378 score cells on the same 168000 target packets per view. CRIC+DURA obtains clean accuracy $85.1429\pm0.4055\%$, six-satellite-view mean $78.1249\pm0.3416\%$, and low-urban accuracy $68.1339\pm0.6603\%$. The corresponding un-restored CRIC values are $85.1139\pm0.3799\%$, $43.9855\pm5.7410\%$, and $33.6470\pm2.9346\%$. Six-view means are formed within each seed before computing the cross-seed sample SD.

All three CRIC seeds share one frozen G, so their SD estimates F-seed variation conditional on G, not variability of three independently trained complete systems. The table compares complete native recipes on a common data split: external supervised/ordinary-PL runs use satellite concatenation, whereas original CRIC uses its native unlabeled mechanism and no As branch. Their difference cannot be attributed solely to cross-RX contrast. CE controls stronger than their Full counterparts remain visible. Additional matched controls below are prospective; no results are inferred for untested data contracts.

\subsubsection{Prospective Coverage and Data-Matched Controls}
The extended design covers F0/F1/F2, chronological D3, ManyRx, label fractions 1/5\%, and ten source-count conditions. The latter use $m\in\{1,2,3,4,6\}$ along the fixed order $(1,3,4,6,8,11)$, under natural or fixed annotation budgets. Their target set is always $(0,2,5,7,9,10)$, six RX and 144000 packets. No new five-source sweep point is trained or queued. Existing F0 five-source results use seven target RX and must not be inserted into this curve.

For natural budgets, $(|L|,|U|,|V|)=(1260m,11340m,5400m)$. The fixed-budget setting uses $(1080,12600m-1080,5400)$, with surplus V excluded. A single source provides no cross-RX positive: contrast is finite zero and effective-anchor coverage is explicitly zero, rather than inventing RX domains from dates or augmentation. Because U still grows with m, the semi-supervised curve does not isolate receiver diversity alone. The strict-supervised fixed-budget curve is the cleaner control for annotation and update count.

Each of the 17 data conditions has native CRIC+DURA, a satellite-concatenated supervised control, and its ordinary-PL control, each with three F seeds. The latter two preserve the same representation and supervised regularizers; PL alone adds current-student hard labels at threshold 0.95, weight one, and U/L ratio one, without EMA or complex gating. They are named PAMF-Sup and PAMF-PL, not relabeled as original CRIC. F training fixes 11100 updates (50 logical epochs of 222 steps), with L batch128. On changed source pools, logical epochs are not claimed to be full dataset passes. Native CRIC retains the 33/17 schedule and U batch256; matched PL starts in the first batch with U batch128. Actual unique observations, repeats, partial batches, and acceptance are reported.

The 153 F instances comprise three exact-contract F0 reuses and 150 planned new models. G-L is trained using L only for a strict-supervised complete system; G-LU uses L plus TX-hidden U. An F without PL attached to G-LU is additionally evaluated, but is described as F-supervised/G-unlabeled, not a supervised system. Each contract initially fixes one G seed shared across F seeds. Two further G seeds at the predetermined three-source natural condition yield a $3\times3$ crossed F/G sensitivity panel. This gives 36 G sources, one existing reuse and 35 planned new complete training lineages. Changing RX, date, TX or physical roles requires matching F and all G ancestors; an F0 teacher cannot initialize another contract. These extensions have not yet been launched.

\subsection{Prospective Mechanism and Efficiency Experiments}
\label{sec:cric-ablation}
All new settings are fixed before additional target predictions and retain clean plus six practical views. Predictions for all views of each row precede independent truth-side scoring. We report Accuracy, Macro-F1, worst-RX accuracy, seed-paired differences, and repaired/newly introduced errors on the same physical observations. Target scores do not select epochs, variants, seeds, or reruns.

\subsubsection{Identification and Restoration Contributions}
The principal F0 factorial is BASE versus CRIC, crossed with raw versus DURA. BASE retains the original R-Drop/SSL recipe but lacks the added contrast; it is not pure CE. Existing BASE, ALL\_TX and CROSS\_RX each have three E50 models. Their two readouts yield 18 frozen model/readout combinations and 126 view cells; missing seven-view predictions must be completed rather than inferred from existing clean results. The interaction is
\begin{equation}
\begin{split}
I={}&(A_{\mathrm{CRIC},G}-A_{\mathrm{CRIC},raw})\\
&-(A_{\mathrm{BASE},G}-A_{\mathrm{BASE},raw}).
\end{split}
\end{equation}
ALL\_TX with the same G tests whether restricting positives to different physical RX improves on ordinary same-TX contrast.

F ablations separately use shallow-only, intermediate-only, and final-embedding-only contrast; zero R-Drop KL with two CE forwards; one CE forward; one SIG\_TIME projection; removal of the time coordinate; removal of SIG\_TIME, PA or reference response; student-generated rather than EMA pseudo labels; removal of neighbor or strong-agreement selection; removal of the U entropy term; and a same-representation CE-only labeled reference. The first fourteen settings preserve unaffected components; the CE-only reference changes the complete training recipe and is not a single-factor effect. These 15 prospective variants use three seeds each, separately from the 153-instance comparison matrix. Structural removals require matched training; frozen branch masking is only a sensitivity diagnostic.

G ablations include coarse-only, priors without CG, no conjugate-image term, and omission of the backprojected-residual prior input. Initialization copying and stage-IQ distillation are crossed in a $2\times2$ design so their effects are not conflated. Four structural and three additional initialization/distillation settings give 21 planned final-student runs over three seeds, with teacher ancestry and cost reported separately. Frozen-weight solver sensitivity uses CG counts 1/2/4 per stage, prior-scale multipliers 0.5/1/2, and a one-stage path. These are distinct runtime variants requiring formal seven-view tests, not evidence for a retrained causal ablation or permission to replace the default according to target rank.

\subsubsection{Latency Boundary and Repetition}
We measure F alone, G alone, and the actual online G+F call. Three boundaries are separated: synchronized GPU-resident calls including input replacement, owned outputs and validity checks; pinned-CPU IQ to CPU class including H2D/D2H; and cold loading/warmup/Graph capture. Cached restored-IQ reads never replace G in a complete-system timing. Module profiles are separate diagnostics; their medians are not added to synthesize end-to-end latency.

The common benchmark uses one exclusive RTX3090, the same Python3.10/Torch2.1 stack, two CPU threads, FP32 and both TF32 flags disabled. Existing Torch1.11 comparator costs are not ranked directly against DURA's Torch2.1 timing. Every method has an eager result and, where supported, a Graph result; unsupported capture is reported as N/A with its reason. Healthy training is not interrupted to obtain a timing GPU.

Batches are 1/8/32/64/128. Five balanced-order rounds use 100 warmups and 1000 measured calls per round for B1, or 20 warmups and 100 measured calls for larger batches. We report synchronized median/P95, sample counts, between-round variation, whole-batch latency, throughput, resident/incremental peak memory, and cold setup. CUDA-event device time is a separate metric. Actual input replacement, output ownership, numerical/argmax agreement and unchanged frozen state are checked without query access. Three model seeds and five timing rounds describe different sources of variability. Amortized batch time is not single-observation waiting latency.

\begin{table}[t]
\centering\footnotesize
\caption{Same-weight DURA execution ablation. T1--T4 isolate the RR/Graph interaction; T0--T1 also changes tensorization and synchronization.}
\label{tab:dura-execution}
\begin{tabular}{lccc}
\toprule
Path & Tensorized checks & RR & Graph\\
\midrule
T0 native & No & Native & No\\
T1 & Yes & No & No\\
T2 & Yes & Yes & No\\
T3 & Yes & No & Yes\\
T4 default & Yes & Yes & Yes\\
\bottomrule
\end{tabular}
\end{table}

Table~\ref{tab:dura-execution} is measured for G and G+F. The earlier seed01 B1 online G+F median 4.531972ms (P95 4.990450ms) is a historical GPU-resident measurement including engine copies/checks, excluding transfer, queueing and first capture. It is not the result of this prospective common-stack benchmark. A 256-sample window at 25MS/s spans 10.24$\mu$s, so this latency alone does not establish continuous full-rate real-time processing.

Deployment counts include the learned frozen classifier even when exported as a buffer, and exclude optimizer, EMA and unused training branches. Convolution/linear MACs, FFT/IFFT counts and complex solver costs are reported separately. R-Drop's second pass, contrast construction, EMA, U gates and entropy are training-only; their per-step forward/backward/optimizer time and memory are measured during training, not assigned fictional inference passes. Contrast reuses F forwards but still incurs descriptor/similarity computation. Total training cost includes all G ancestors and teachers, unique/repeated L/U exposures, simulation, prediction and failed attempts. Pending timings and ablation scores remain unreported until measured.
